\chapter{2つの銀河を接近させる}
\label{chap:interaction}

第4章では、一つの固定したPlummerポテンシャル中に、回転するテスト粒子円盤を作りました。この章では、その円盤を二つ用意して互いに接近させます。第2章の重心と相対運動、第3章のPlummer重力、第4章の円盤生成を組み合わせ、本書の中心となる銀河相互作用の計算を完成させます。

新しい数値積分法は導入しません。時間発展には、引き続き加速度を再利用するKDK leapfrogを使います。今回新しく考えるのは、\textbf{銀河中心も星も動く状況で、どの位置から、どの粒子への重力を計算するか}です。初期状態を作り、中心軌道を確認してから、潮汐テイルやブリッジを見ます。見た目の変化と数値誤差を区別するため、同じ初期粒子を使った時間刻みの比較も行います。

\begin{goalbox}[この章で学ぶこと]
\begin{itemize}
\item 移動するPlummerポテンシャルと、星を表すテスト粒子の役割を区別できます。
\item 採用した中心間ポテンシャルに合わせて、近点距離を指定した初期軌道を作れます。
\item 円盤の局所的な位置・速度を慣性系へ移し、銀河中心と星を同じ時刻へ更新できます。
\item 中心軌道の保存量と、時間に依存する重力場中の粒子のエネルギー変化を区別できます。
\item 潮汐構造を図と粒子割合で調べ、時間刻みによる結果の違いを確認できます。
\end{itemize}
\end{goalbox}

\section{これまでの計算要素を一つにまとめる}
\label{sec:ch5-model}
\subsection{銀河を変形しない重力場として表す}
各銀河$k=1,2$に、一定の全質量$M_k$とPlummerスケール長$a_k$を与えます。銀河中心の位置・速度を$\Rvec_k,\Vvec_k$とし、その中心を動かします。一方、$M_k$と$a_k$自体は時間とともに変えません。したがって、移動するのは\textbf{形を固定した解析的ポテンシャル}です。Plummer球を多数の質量粒子で表し、その変形まで計算するわけではありません。

円盤の星は、質量ゼロのテスト粒子$\rvec_i,\vvec_i$です。星は両銀河の重力を受けますが、星同士の重力や、星から銀河中心への反作用は計算しません。粒子の表示上の重みを、Plummer球の質量と足し合わせることもしません。第4章の粒子数面密度$\Sigma_N$と、重力を作る質量分布の区別を引き継ぎます。

\begin{table}[htbp]
\centering
\caption{本章で計算する重力の関係。「力を受ける側」と「重力を作る側」を区別します。星には出身銀河だけでなく、両銀河の重力を与えます。}
\label{tab:ch5-coupling}
\begin{tabularx}{\linewidth}{lX}
\toprule
力を受ける側 & 計算する重力 \\
\midrule
銀河中心1・2 & 中心間の有効対ポテンシャルから求める相互力 \\
テスト粒子 & 二つの移動するPlummerポテンシャルの重力の和 \\
テスト粒子間 & 計算しません \\
テスト粒子から中心 & 計算しません \\
\bottomrule
\end{tabularx}
\end{table}

外部重力、ガス、星形成、動力学的摩擦は含めません。すべての位置と速度を一つの慣性系で計算し、運動は$x$--$y$平面内に限定します。円盤が変形しても、背景のPlummer球は同じ形です。ここで再現するのは、簡略化した重力モデルに対する星の軌道の応答です。

\subsection{本章で統合する要素を確認する}
第2章の相対位置・相対速度を
\begin{equation}
\Dvec=\Rvec_2-\Rvec_1,\qquad D=|\Dvec|,\qquad \Uvec=\Vvec_2-\Vvec_1
\label{eq:ch5-relative}
\end{equation}
と書きます。銀河中心の軌道を決めた後、第4章の円盤を各中心の周りへ配置します。必要な操作は、重心系の初期軌道、円盤の初期位置・速度、加速度の計算、KDKの時間更新の四つです。

台形則と線形補間は円盤のCDF生成で再利用します。求根法は第3章の半質量半径などを求める道具として既習ですが、\textbf{本章の標準計算では新たな求根を無理に挿入しません}。後述する中心軌道の初期速度は、保存量から直接求められるためです。各手法は、必要な役割を持つ場所で使います。

\section{動いている中心から重力を求める}
\label{sec:ch5-gravity}
\subsection{テスト粒子が受ける二つのPlummer場を合成する}
銀河$k$の中心から星までの相対位置を
\begin{equation}
\boldsymbol s_{k,i}=\rvec_i-\Rvec_k
\label{eq:ch5-local-position}
\end{equation}
とします。第3・4章の原点からの位置を、この相対位置へ置き換えます。銀河$k$によるポテンシャルと加速度は
\begin{align}
\Phi_k(\rvec_i,t)&=-\frac{GM_k}{\sqrt{|\boldsymbol s_{k,i}|^2+a_k^2}},
\label{eq:ch5-moving-potential}\\
\avec_{k,i}&=-GM_k\frac{\boldsymbol s_{k,i}}{(|\boldsymbol s_{k,i}|^2+a_k^2)^{3/2}}
\label{eq:ch5-moving-acceleration}
\end{align}
です。Plummer球の形は固定でも、中心が動くため、慣性系ではポテンシャルが時間に依存します。Plummer場の式そのものは前章と同じです。\cite{bt2008}

まず1粒子の$x$成分だけを見ると、中心位置$X_k$を使って
\begin{equation}
a_{x,k,i}=-\frac{GM_k(x_i-X_k)}{[(x_i-X_k)^2+(y_i-Y_k)^2+a_k^2]^{3/2}}
\label{eq:ch5-component}
\end{equation}
です。原点ではなく、動く銀河中心へ向かう加速度であることを確認してください。前章の配列計算では、最初に\texttt{pos - center}を計算するだけで同じ置き換えができます。
\inputcode{chap05/moving-plummer}

各星の運動方程式は
\begin{equation}
\frac{d\rvec_i}{dt}=\vvec_i,\qquad \frac{d\vvec_i}{dt}=\avec_{1,i}+\avec_{2,i}
\label{eq:ch5-tracer-ode}
\end{equation}
です。どちらの円盤から出発した星も、同じ二つの場の和を受けます。出身銀河の番号は色分けと解析に使うだけで、重力源を選別する番号ではありません。

\subsection{中心間の相互作用を別に定義する}
中心の運動にも、有限の大きさを反映した滑らかな力を使います。本章では二つの中心の有効な対ポテンシャルエネルギーを
\begin{equation}
W_{\rm c}(D)=-\frac{GM_1M_2}{\sqrt{D^2+b^2}}
\label{eq:ch5-pair-potential}
\end{equation}
と\textbf{モデルとして定義}します。$b$は中心間力に使う共通の有効長さです。$a_1,a_2$は星が受ける各Plummer場のスケール長なので、記号を分けます。

式\eqref{eq:ch5-pair-potential}は、\textbf{二つの広がったPlummer密度分布どうしの厳密な相互作用エネルギーではありません}。その計算には、各分布の広がりにわたって相互作用を足す必要があります。本章はその積分を追加せず、中心軌道のための簡単な有効モデルを採用します。後で与える$b=3\unit{kpc}$も、この教材の選択です。

同じ$W_{\rm c}$から両方の力を求めると、中心加速度は
\begin{equation}
\boldsymbol A_1=+GM_2\frac{\Dvec}{(D^2+b^2)^{3/2}},\qquad
\boldsymbol A_2=-GM_1\frac{\Dvec}{(D^2+b^2)^{3/2}}
\label{eq:ch5-center-acceleration}
\end{equation}
になります。この二つは$M_1\boldsymbol A_1+M_2\boldsymbol A_2=0$を満たします。第2章と同じく、共通部分を一度だけ計算します。
\inputcode{chap05/center-acceleration}

$b$は両方の式で同じ値です。$M_1\ne M_2$の場合でも作用・反作用を保ちます。$a_1\ne a_2$へ拡張するとき、中心1へは$a_2$、中心2へは$a_1$をそのまま代入すると、一般に二つの力の大きさが一致しません。その置き換えは行いません。本章の主実験では$a_1=a_2=b$に固定します。

この選択には、力の対応を単純にする利点があります。$b=a_2$なら、中心1への加速度は銀河2のPlummer場を$\boldsymbol R_1$で評価した値に一致します。中心2についても同じ対応を保つため、標準では$b=a_1=a_2$とします。その結果、星と出身中心が受ける相手の重力の差を、同じ場の二つの位置での評価として書けます。ただし、これは本章の有効モデル内での整合性であり、広がった二つの球の厳密な相互力を得たという意味ではありません。


\begin{failbox}{Plummer場と中心間力を同じ意味にしない}
星は式\eqref{eq:ch5-moving-potential}の場を受け、中心は式\eqref{eq:ch5-pair-potential}で動きます。形が似ていても、後者は本章で定義した有効な粒子対のモデルです。有限サイズの銀河を厳密に剛体として相互作用させた計算、と説明しないようにします。
\end{failbox}


% BEGIN TikZ addition: ch05_force_model
図\ref{fig:concept-ch5-forces}に、このモデルで計算する作用をまとめました。中心間の有効長さ$b$と、星が受けるPlummer場の$a_k$を区別し、矢印のない反作用を計算へ加えていないことを確認してください。
\begin{figure}[!htbp]
\centering
% Dependency arrows, not spatial force vectors.
\begin{tikzpicture}[textbook diagram]
  \node[anchor=west,font=\sffamily\bfseries] at (0,4.45) {Which objects affect which others?};
  \node[draw=figblue,fill=figblue!5,rounded corners=1pt,text width=33mm,minimum height=16mm,align=center] (c1) at (2.0,2.80) {Center 1\\$M_1,\ \boldsymbol R_1$};
  \node[draw=figred,fill=figred!5,rounded corners=1pt,text width=33mm,minimum height=16mm,align=center] (c2) at (9.7,2.80) {Center 2\\$M_2,\ \boldsymbol R_2$};
  \draw[{Latex}-{Latex},line width=0.85pt] (c1.east) -- node[above,align=center] {Effective pair\\potential: $b$} (c2.west);
  \node[draw=black!55,fill=white,rounded corners=1pt,text width=103mm,minimum height=17mm,align=center] (stars) at (5.85,-0.90) {Massless tracers: $\boldsymbol r_i$\\Every tracer feels both Plummer fields.};
  \draw[vector,figblue] (c1.south) -- node[right,align=left,pos=.52] {Plummer field 1\\$M_1,\ a_1$} (2.0,-.04);
  \draw[vector,figred,dashed] (c2.south) -- node[left,align=right,pos=.52] {Plummer field 2\\$M_2,\ a_2$} (9.7,-.04);
  \node[draw=black!30,rounded corners=1pt,text width=103mm,align=center,inner sep=5pt] at (5.85,-2.64) {No tracer--tracer forces.\\No feedback from tracers to the centers.};
  \node[small label,text=figgray,anchor=west] at (0,-3.65) {Arrows show model dependencies, not force directions in space.};
\end{tikzpicture}

\caption{中心とテスト粒子の間で計算する作用の関係。横の双方向矢印は共通の有効長さ$b$を使う中心間相互作用、下向き矢印は各Plummer場から星への作用です。星から中心への反作用と星同士の重力は計算しません。矢印はモデルの依存関係であり、空間内の力の向きではありません。}
\label{fig:concept-ch5-forces}
\end{figure}

% END TikZ addition: ch05_force_model

\Needspace{10\baselineskip}
\section{近点距離に合った初期軌道を作る}
\label{sec:ch5-initial-orbit}
\subsection{点質量の速度をそのまま使わない}
銀河中心の相対加速度は、式\eqref{eq:ch5-center-acceleration}の差から
\begin{equation}
\frac{d\Uvec}{dt}=-\frac{\mu\Dvec}{(D^2+b^2)^{3/2}},\qquad \mu=G(M_1+M_2)
\label{eq:ch5-relative-ode}
\end{equation}
です。その相対運動の比エネルギーと比角運動量は
\begin{equation}
\varepsilon=\frac{U^2}{2}-\frac{\mu}{\sqrt{D^2+b^2}},\qquad h=(\Dvec\times\Uvec)_z
\label{eq:ch5-orbital-invariants}
\end{equation}
で、どちらも時間によらず一定になります。第2章の$-\mu/D$とはポテンシャルが異なります。同じ$\mu$と近点距離を選んでも、初期速度は改めて求めます。

本章では$\varepsilon=0$の内向きの遭遇を作ります。これは無限遠をポテンシャルの零点とした、束縛・非束縛の境界のエネルギーです。ただし、$b>0$では力が厳密な逆2乗則ではないため、\textbf{軌道の形をKepler放物線と呼ぶことはできません}。$b\to0$の極限で、第2章の放物線に対応します。

\subsection{近点から角運動量を決める}
初期距離を$D_0$、目標の近点距離を$D_p$とし、$0<D_p<D_0$を選びます。近点では動径速度が0なので、速度の大きさは接線速度$U_p$です。$\varepsilon=0$を使うと
\begin{equation}
U_p^2=\frac{2\mu}{\sqrt{D_p^2+b^2}},\qquad
h^2=D_p^2U_p^2=\frac{2\mu D_p^2}{\sqrt{D_p^2+b^2}}
\label{eq:ch5-pericenter-angular}
\end{equation}
です。ここでは$h>0$を選び、軌道の向きを反時計回りにします。

初期位置を$\Dvec(0)=(D_0,0)$とすれば、角運動量から接線速度$U_t=h/D_0$が決まります。初期位置でも$\varepsilon=0$とし、動径速度の内向きの符号を選ぶと
\begin{equation}
U_t=\frac{h}{D_0},\qquad
U_r=-\sqrt{\frac{2\mu}{\sqrt{D_0^2+b^2}}-U_t^2},\qquad
\Uvec(0)=(U_r,U_t)
\label{eq:ch5-incoming-speed}
\end{equation}
となります。これは採用した有効ポテンシャルから直接求めた初期条件で、近点を数値計算後に合わせ直す操作ではありません。

\subsection{重心系の二つの位置・速度へ変換する}
第2章の重心変換を再利用します。初期重心を原点、重心速度を0とすれば、$\Mtot=M_1+M_2$として
\begin{align}
\Rvec_1(0)&=-\frac{M_2}{\Mtot}\Dvec(0),&
\Rvec_2(0)&=+\frac{M_1}{\Mtot}\Dvec(0),\notag\\
\Vvec_1(0)&=-\frac{M_2}{\Mtot}\Uvec(0),&
\Vvec_2(0)&=+\frac{M_1}{\Mtot}\Uvec(0)
\label{eq:ch5-barycentric-initial}
\end{align}
です。相対距離・相対速度を、それぞれの中心から原点までの距離・速度と取り違えないようにします。
\inputcode{chap05/encounter-initial}

返り値は形$(2,4)$の配列です。各行に銀河中心の$[X_k,Y_k,V_{x,k},V_{y,k}]$を入れます。\texttt{weights}は式\eqref{eq:ch5-barycentric-initial}の二つの係数です。\texttt{[:, None]}で列を加え、各係数を$x,y$の両成分に掛けています。

\section{二つの円盤を配置する}
\label{sec:ch5-disk-initial}
\subsection{円盤の内部運動と銀河全体の運動を分ける}
各円盤の粒子数面密度は、第4章と同じ切断指数分布です。
\begin{equation}
\Sigma_{N,k}(R)\propto e^{-R/R_{d,k}},\qquad 0\le R\le R_{\max,k}.
\label{eq:ch5-disk-profile}
\end{equation}
この$R$は各銀河中心からの円盤半径です。全体系の原点から測る半径ではありません。CDF表を台形則で作り、その逆関係を線形補間で読む処理は、前章のものを再利用します。

円盤の局所座標を$\boldsymbol s_i$、局所速度を$\boldsymbol w_i$と書きます。粒子の半径$R$での円運動速度は
\begin{equation}
v_{{\rm c},k}^2(R)=\frac{GM_kR^2}{(R^2+a_k^2)^{3/2}}
\label{eq:ch5-circular-speed}
\end{equation}
です。これに局所の接線方向を掛け、各銀河の周りを回る速度を作ります。\inputcode{chap05/make-disk}

慣性系へ移すときには、位置だけでなく速度にも銀河中心の値を足します。
\begin{equation}
\rvec_i(0)=\Rvec_k(0)+\boldsymbol s_i,\qquad
\vvec_i(0)=\Vvec_k(0)+\boldsymbol w_i.
\label{eq:ch5-place-disk}
\end{equation}
\inputcode{chap05/place-disk}

この変換は、円盤粒子だけの平均位置を引く「再中心化」とは違います。第4章の各粒子とPlummer中心の相対位置・速度をそのまま保ち、銀河全体へ並進運動を与えています。


% BEGIN TikZ addition: ch05_frame_transform
図\ref{fig:concept-ch5-frame}の左右を対応させると、位置と速度にそれぞれ何を足すかが分かります。局所速度$\vect w_i$だけでは、銀河全体の並進運動は含まれません。
\begin{figure}[!htbp]
\centering
% Separate diagrams for position addition and velocity addition; lengths have different units.
\begin{tikzpicture}[textbook diagram]
  \node[anchor=west,font=\sffamily\bfseries] at (0,4.80) {Add the center position and the center velocity};
  \node at (2.75,4.18) {Position space};
  \node at (9.00,4.18) {Velocity space};
  \coordinate (O) at (0.45,0.68); \coordinate (C) at (2.93,1.50); \coordinate (P) at (4.55,3.32);
  \coordinate (Z) at (6.60,0.68); \coordinate (V) at (9.32,1.40); \coordinate (W) at (10.85,3.30);
  \draw[vector,figblue] (O) -- node[below,pos=.52] {$\boldsymbol R_k$} (C);
  \draw[vector,figred,dashed] (C) -- node[right,pos=.60] {$\boldsymbol s_i$} (P);
  \draw[vector] (O) -- node[above left,pos=.56] {$\boldsymbol r_i$} (P);
  \draw[vector,figblue] (Z) -- node[below,pos=.52] {$\boldsymbol V_k$} (V);
  \draw[vector,figred,dashed] (V) -- node[right,pos=.57] {$\boldsymbol w_i$} (W);
  \draw[vector] (Z) -- node[above left,pos=.53] {$\boldsymbol v_i$} (W);
  \fill (O) circle(1.3pt);\node[below left] at (O) {$O$};
  \fill (Z) circle(1.3pt);\node[below left] at (Z) {$0$};
  \draw[black] ($(C)+(-.07,-.07)$) -- ($(C)+(.07,.07)$);\draw[black] ($(C)+(-.07,.07)$) -- ($(C)+(.07,-.07)$);
  \fill (P) circle(1.8pt);\node[above right] at (P) {Star $i$};
  \node[small label,below,yshift=-1.0mm] at (C) {Center $k$};
  \draw[black!20] (5.78,0.3) -- (5.78,3.9);
  \node at (2.75,-0.13) {$\boldsymbol r_i=\boldsymbol R_k+\boldsymbol s_i$};
  \node at (9.00,-0.13) {$\boldsymbol v_i=\boldsymbol V_k+\boldsymbol w_i$};
  \node[small label,text=figgray] at (5.80,-0.88) {Blue: center motion; dashed red: motion relative to the center.};
\end{tikzpicture}

\caption{局所円盤を慣性系へ移す二つのベクトル加算。位置には中心位置$\Rvec_k$、速度には中心速度$\Vvec_k$を加えます。朱色破線は中心に対する局所量です。左右は位置空間と速度空間の図なので、矢印の長さを相互に比較しません。}
\label{fig:concept-ch5-frame}
\end{figure}
\FloatBarrier
% END TikZ addition: ch05_frame_transform

\begin{failbox}{中心速度も足します}
位置にだけ$\Rvec_k$を足して、速度には$\Vvec_k$を足さないと、円盤が銀河中心と一緒に進みません。その後の崩れを潮汐相互作用の結果と判断してはいけません。初期値から$\rvec_i-\Rvec_k$と$\vvec_i-\Vvec_k$を計算し、局所円盤へ戻るか確認します。
\end{failbox}

\subsection{progradeの意味を定める}
第4章では円盤の向きを反時計回り・時計回りと呼びました。今回は二つの銀河の軌道があるので、円盤の角運動量と軌道の角運動量を比較できます。同じ向きの回転を\textbf{prograde}、逆向きを\textbf{retrograde}と呼びます。

標準軌道は$h>0$で反時計回りです。したがって、両円盤で\texttt{spin=+1}を選べばprograde--progradeの遭遇になります。$x$軸の右側・左側のどちらに銀河があるかではなく、\textbf{中心に対する局所的な回転の向き}で判断します。本章はまずこの一例を完成させ、向きを変える比較は章末の演習と第6章へ残します。


% BEGIN TikZ addition: ch05_spin_sense
図\ref{fig:concept-ch5-spin}は、軌道の向きを固定したまま円盤2の内部回転だけを変える比較です。中心の軌道方向を示す大きな矢印と、各円盤の回転を示す小さな矢印を区別してください。
\begin{figure}[!htbp]
\centering
% Sense-of-rotation sketch. The large curved arrows do NOT depict the center trajectories.
\begin{tikzpicture}[textbook diagram]
  \node[anchor=west,font=\sffamily\bfseries] at (0,4.70) {Change the internal spin, not the center orbit};
  \foreach \shift/\title/\mode in {0/Prograde--prograde/1,6.1/Prograde--retrograde/-1} {
    \begin{scope}[shift={(\shift,0)}]
      \node[font=\sffamily\bfseries] at (2.65,4.05) {\title};
      \node[small label] at (2.65,3.60) {Same orbital sense: $h>0$};
      \draw[vector,black] (3.64,2.57) arc[start angle=28,end angle=148,radius=1.18];
      \node[small label,fill=white,inner sep=1pt] at (2.65,2.71) {Orbit};
      \foreach \cx/\col/\id in {1.18/figblue/1,4.12/figred/2} {
        \draw[\col,densely dashed] (\cx,1.55) circle(.86);
        \draw[black] ({\cx-.07},1.48) -- ({\cx+.07},1.62);\draw[black] ({\cx-.07},1.62) -- ({\cx+.07},1.48);
        \node[below] at (\cx,.52) {Disk \id};
      }
      \draw[vector,figblue] (1.18,1.55) ++(-45:.56) arc[start angle=-45,end angle=238,radius=.56];
      \ifnum\mode=1
        \draw[vector,figred,dashed] (4.12,1.55) ++(-45:.56) arc[start angle=-45,end angle=238,radius=.56];
        \node at (4.12,-.10) {$s_{{\rm rot},2}=+1$};
      \else
        \draw[vector,figred,dashed] (4.12,1.55) ++(238:.56) arc[start angle=238,end angle=-45,radius=.56];
        \node at (4.12,-.10) {$s_{{\rm rot},2}=-1$};
      \fi
      \node at (1.18,-.10) {$s_{{\rm rot},1}=+1$};
    \end{scope}
  }
  \draw[black!20] (5.73,-.40) -- (5.73,4.2);
  \node[small label,text=figgray] at (5.65,-.93) {Curved arrows indicate rotation senses only, not orbit shapes.};
\end{tikzpicture}

\caption{軌道の向きを$h>0$に固定した回転方向の比較。左は両円盤がprograde、右は円盤2の内部回転だけがretrogradeです。大きな曲線矢印は軌道の回転の向きを示す記号で、中心軌道の形ではありません。回転は銀河中心に対する局所的な向きで判断します。}
\label{fig:concept-ch5-spin}
\end{figure}
\FloatBarrier
% END TikZ addition: ch05_spin_sense

\subsection{二銀河系の初期配置の近似を確認する}
局所の円運動速度は、出身銀河だけのPlummer場に対して決めています。有限の$D_0$から計算を始めるので、相手の銀河の重力は時刻0から加わります。したがって、これは孤立円盤二つを初期条件として配置したものであり、二銀河の合成場に対する厳密な平衡円盤ではありません。

初期距離を十分離して選ぶのは、この初期の摂動を抑えるためでもあります。ただし、相手の力を初めの数ステップだけ無視したり、説明なく徐々に強めたりはしません。式\eqref{eq:ch5-tracer-ode}の力を最初から一貫して使います。

\section{全体の配列を作り、同じ時刻へ進める}
\label{sec:ch5-integration}
\subsection{配列内の各行の役割を決める}
テスト粒子の合計を$N=N_1+N_2$とします。\texttt{state}は形$(2+N,4)$で、最初の2行を銀河中心、その後を円盤1と円盤2の星にします。
\begin{center}
\begin{tabular}{lll}
\toprule
Pythonの行 & 対象 & 列 \\
\midrule
\texttt{state[0]} & 銀河中心1 & \texttt{[x, y, vx, vy]} \\
\texttt{state[1]} & 銀河中心2 & 同上 \\
\texttt{state[2:2+N1]} & 円盤1のテスト粒子 & 同上 \\
\texttt{state[2+N1:]} & 円盤2のテスト粒子 & 同上 \\
\bottomrule
\end{tabular}
\end{center}

星にだけ対応する長さ$N$の整数配列\texttt{origin}も作ります。0が銀河1、1が銀河2の出身です。Pythonの番号は0から始まるので、本文の銀河番号より1小さい点に注意します。この出身番号は時間とともに変えません。他方の銀河の近くへ移った粒子も、同じ色と番号で追跡します。
\inputcode{chap05/assemble-state}

乱数生成器は一つだけ作り、円盤1、円盤2の順に使います。円盤ごとに同じseedへ戻して、全く同じ点配置を二つ作ることはしません。以後の時間刻み試験では、ここで得た初期配列を保存し、その同じ配列を使います。

\subsection{一度の呼び出しで全加速度を返す}
中心への加速度と星への加速度を、一つの関数へまとめます。ただし、その計算内容は分けたままです。
\inputcode{chap05/system-acceleration}

返り値は入力の\texttt{pos}と同じ形$(2+N,2)$です。先頭2行は式\eqref{eq:ch5-center-acceleration}、残りは式\eqref{eq:ch5-tracer-ode}です。後半の星の計算に\texttt{origin}は使いません。星の出身にかかわらず、二つの重力を足すためです。

中心同士の計算は1組だけです。星には二つの解析的な場を評価するだけなので、1回の全加速度計算は$O(N)$です。$N\times N$の距離表や、星の粒子対を走査する二重ループは作りません。

\subsection{KDKの更新式をそのまま使う}
第4章のKDKに、今作った加速度関数を渡します。
\inputcode{chap05/leapfrog-step}

最初の半kickで全行の速度を更新し、driftで全行の位置を更新します。その\texttt{pos\_new}には、星だけでなく\textbf{新しい銀河中心の位置}も入っています。そのため、次の加速度は新しい中心と新しい星の位置から計算されます。最後の半kickを終えると、全行の位置と速度が同じ整数時刻にそろいます。

\begin{failbox}{動く中心と星の時刻をずらさない}
中心だけを先に1ステップ進め、星には古い中心位置を使う処理では、ここで定めたKDKになりません。全行のdriftを終えてから、一つの位置配列を使って次の重力を求めます。描画用に間引いた中心の軌道を、時間積分の途中で代用することもしません。
\end{failbox}

力が位置から決まるので、この全行の更新を使えます。ただし、反作用を省いたテスト粒子を含む全配列を、通常の自己重力N体系と同一視しません。中心の二体系には保存量がありますが、星を加えた表示用のエネルギーの総和を、そのまま保存量にはできません。実際に確かめる量は\ref{sec:ch5-center-check}節と\ref{sec:ch5-energy}節で分けます。

\section{標準モデルを計算する}
\label{sec:ch5-standard}
\subsection{パラメータをそろえる}
本章の標準条件を表\ref{tab:ch5-standard}に示します。質量、Plummerスケール長、円盤分布は既存の標準モデルを引き継ぎます。二つの銀河を初期距離60 kpcから、目標近点8 kpcへ向かわせます。これらは教材用の値で、特定の観測銀河を再現するための調整ではありません。

\begin{table}[htbp]
\centering
\caption{第5章の標準モデル。長さの役割を区別します。全質量は$M_1+M_2=2\times10^{11}\Msun$で、テスト粒子に質量を加えません。}
\label{tab:ch5-standard}
\begin{tabular}{ll}
\toprule
量 & 設定 \\
\midrule
重力源 & $M_1=M_2=10^{11}\Msun$ \\
各Plummer場 & $a_1=a_2=3\unit{kpc}$ \\
中心間有効長さ & $b=3\unit{kpc}$ \\
円盤初期分布 & 両方$R_d=3\unit{kpc}$、$R_{\max}=15\unit{kpc}$ \\
粒子数・乱数 & $N_1=N_2=1000$、seed = 42 \\
CDF表 & 各円盤4096区間、4097分点 \\
中心軌道 & $D_0=60\unit{kpc}$、$D_p=8\unit{kpc}$、$\varepsilon=0$ \\
回転方向 & 軌道も両円盤も反時計回り \\
積分時間・時間刻み & $t_{\rm end}=0.60\unit{Gyr}$、$\Delta t=\chfivevalDt\unit{Myr}$ \\
ステップ数・保存 & \chfivevalNstep ステップ、星は\chfivevalSaveEvery ステップごと \\
\bottomrule
\end{tabular}
\end{table}

第4章の孤立円盤では$\Delta t=0.2\unit{Myr}$を使いました。本章では近接遭遇の後の粒子軌道も追うため、\textbf{時間刻みを$0.05\unit{Myr}$へ選び直します}。この選択は孤立円盤の誤差をそのまま流用したものではなく、\ref{sec:ch5-convergence}節の実際の刻み試験に基づきます。出力の間隔は2 Myrに保ちます。
\inputcode{chap05/parameters}

初期相対速度の成分は、式\eqref{eq:ch5-incoming-speed}から
\begin{equation}
U_r=\chfivevalVr\unit{km\,s^{-1}},\qquad U_t=\chfivevalVt\unit{km\,s^{-1}}
\label{eq:ch5-initial-values}
\end{equation}
となり、大きさは$\chfivevalSpeed\unit{km\,s^{-1}}$です。等質量なので各中心の初期位置は$(-30,0)$ kpcと$(30,0)$ kpcです。速度はそれぞれ相対速度の$-1/2$倍、$+1/2$倍になります。コード内部では速度をkpc/Gyrへ換算しています。

\subsection{中心と星で保存間隔を分ける}
近点時刻や保存量を調べるため、銀河中心の2行だけは毎ステップ保存します。星を含む全配列は40ステップごとに保存します。これは時間積分を分けるという意味ではありません。\textbf{計算は全行で同じ12000ステップ}を実行し、履歴へ書き込む頻度だけを分けます。
\inputcode{chap05/integrate-leapfrog}

\texttt{prepare\_history}は前章と同じ役割で、初期配列のコピー、保存ステップ番号、時間配列、保存領域を作ります。今回の初期配列は形$(2+N,4)$です。最後の時刻が保存間隔に一致しない場合も、必ず最終状態を保存します。補助的な入力確認のコードはNotebookに収録しています。

標準設定では、全履歴は$(301,2002,4)$、中心の履歴は$(12001,2,4)$です。\texttt{float64}の1要素を8 byteとして、配列だけの容量はそれぞれ\chfivevalHistoryMB\ MB、\chfivevalCenterMB\ MBです。MBは$10^6$ byteで数えています。一時配列、比較計算、図や動画の容量は別です。

全加速度関数の呼び出しは、初回1回と各ステップ1回で\textbf{\chfivevalForce 回}です。1回の呼び出しで、両中心と全テスト粒子の加速度を返します。粒子ごとの場の評価を数えた回数とは区別してください。実行時間はNotebookで測り、特定のノートPCでも同じ秒数になるとは仮定しません。

\section{まず銀河中心の軌道を確かめる}
\label{sec:ch5-center-check}
\subsection{中心二体系の保存量を求める}
テスト粒子からの反作用がないので、中心軌道は星の位置に依存しません。その二体系の全エネルギーは
\begin{equation}
E_{\rm c}=\frac12M_1V_1^2+\frac12M_2V_2^2-\frac{GM_1M_2}{\sqrt{D^2+b^2}}
\label{eq:ch5-center-energy}
\end{equation}
です。中心間のポテンシャルを1回だけ数えます。式\eqref{eq:ch5-pair-potential}で動かした中心に、点質量の$-GM_1M_2/D$を使ってエネルギーを診断しないようにします。

全運動量、原点の周りの全角運動量、重心位置は第2章と同じ
\begin{align}
\Pvec_{\rm c}&=M_1\Vvec_1+M_2\Vvec_2,\notag\\
L_{z,{\rm c}}&=M_1(\Rvec_1\times\Vvec_1)_z+M_2(\Rvec_2\times\Vvec_2)_z,\notag\\
\Rvec_{\rm cm}&=\frac{M_1\Rvec_1+M_2\Rvec_2}{\Mtot}
\label{eq:ch5-center-momentum}
\end{align}
です。$W_{\rm c}$は時間に陽には依存しない中心間のポテンシャルなので、$E_{\rm c}$は一定です。作用・反作用と中心力の関係から、$\Pvec_{\rm c}$と$L_{z,{\rm c}}$も保存され、重心は等速直線運動します。
\inputcode{chap05/center-diagnostics}

\subsection{初期値が0の量で割らない}
重心系の初期エネルギーは理論的には$E_{\rm c}(0)=0$、全運動量も0です。そこで、共通の正の規格化尺度を
\begin{align}
E_*&=\frac{GM_1M_2}{\sqrt{D_0^2+b^2}},&
V_*&=\sqrt{\frac{G\Mtot}{D_0}},\notag\\
P_*&=\Mtot V_*,&
L_*&=\mred D_0V_*,\qquad \mred=\frac{M_1M_2}{\Mtot}
\label{eq:ch5-error-scales}
\end{align}
とします。$V_*$は単位をそろえるための速度の尺度で、有効ポテンシャル中の円運動速度という意味ではありません。

各時刻での初期値からの差の最大値を、$E_*,L_*,P_*$でそれぞれ割り、$e_{E,{\rm c}},e_{L,{\rm c}},e_{P,{\rm c}}$とします。重心の誤差は
\begin{equation}
e_{\rm cm}=\max_n\frac{|\Rvec_{\rm cm}^n-[\Rvec_{\rm cm}(0)+t_n\Pvec_{\rm c}(0)/\Mtot]|}{D_0}
\label{eq:ch5-cm-error}
\end{equation}
です。最大値は、中心を保存した全ステップについて取ります。角運動量の誤差も今回は$L_*$で規格化するため、初期角運動量に対する相対誤差とは呼びません。

\begin{table}[htbp]
\centering
\caption{標準計算での銀河中心の保存量・重心のずれ。規格化は式\eqref{eq:ch5-error-scales}と式\eqref{eq:ch5-cm-error}です。丸め誤差に近い数値の末尾は環境により変わります。}
\label{tab:ch5-center-errors}
\begin{tabular}{lr}\toprule
指標 & 実測値 \\
\midrule
エネルギー $e_{E,{\rm c}}$ & $3.30\times 10^{-6}$ \\
角運動量 $e_{L,{\rm c}}$ & $5.19\times 10^{-15}$ \\
運動量 $e_{P,{\rm c}}$ & $0$ \\
重心位置 $e_{\rm cm}$ & $0$ \\
\bottomrule\end{tabular}

\end{table}

KDKのkickは一組の中心力から両方の速度を変えるため、全運動量と全角運動量を丸め誤差の範囲で保ちます。一方、エネルギーを厳密に保存する方法ではなく、$e_{E,{\rm c}}=\chfivevalEnergyError$のずれが残ります。表示される0は、この等質量・対称な初期条件と実装で得た結果であり、任意の条件でも厳密に0になるという意味ではありません。

\subsection{指定した近点への接近を確かめる}
図\ref{fig:ch5-separation}は銀河中心間距離の時間変化です。標準計算で保存された最小距離は$\chfivevalPeri\unit{kpc}$、その時刻は$\chfivevalTperi\unit{Gyr}$でした。目標の8 kpcへ近づいた後、両中心は離れていきます。
\begin{figure}[htbp]
\centering
\includegraphics[width=0.76\linewidth]{\figpath{ch05_separation}}
\caption{銀河中心間距離。破線は目標の近点8 kpc、四角は毎ステップの保存値の最小です。有限刻みの積分誤差と、真の近点時刻が保存時刻の間にある影響をともに含みます。}
\label{fig:ch5-separation}
\end{figure}

\begin{figure}[htbp]
\centering
\includegraphics[width=0.76\linewidth]{\figpath{ch05_center_energy}}
\caption{中心二体系の符号付きエネルギー差。近点付近で目立つ変化はKDKの有限刻みの効果です。動力学的摩擦を入れた軌道エネルギー損失ではありません。}
\label{fig:ch5-center-energy}
\end{figure}

中心の2行だけを取り出して、テスト粒子を全く置かずに再計算する試験も行います。今回の実装では中心の履歴は全要素で一致しました。星が多数あっても、中心へ反作用を返さないモデルになっていることを確認できます。

\FloatBarrier
\section{潮汐テイルとブリッジを見る}
本書の標準遭遇は、遠方の相手から弱い潮汐摂動を受ける場合ではなく、初期円盤の広がりが深く重なる強い近接遭遇です。目標近点距離$D_{\rm p}=8\unit{kpc}$は、各円盤の初期外縁$R_{\max}=15\unit{kpc}$より小さく、相手の中心が初期円盤の範囲内まで近づきます。この幾何学的な条件だけから、永久に失われる粒子割合は決まりません。以下では遠方展開ではなく、両Plummer場による運動を直接計算します。

\label{sec:ch5-tides}
\subsection{初期状態から近点通過後までを並べる}
図\ref{fig:ch5-snapshots}に四つの時刻を示します。青の丸は銀河1、朱色の三角は銀河2の出身粒子です。黒の$\times$と$+$がそれぞれのPlummer中心です。中心の印は重力場の位置を示すもので、Plummer場に追加した点質量ではありません。

\begin{figure}[htbp]
\centering
\includegraphics[width=0.49\linewidth]{\figpath{ch05_snapshot_000}}\hfill
\includegraphics[width=0.49\linewidth]{\figpath{ch05_snapshot_027}}\\[3mm]
\includegraphics[width=0.49\linewidth]{\figpath{ch05_snapshot_035}}\hfill
\includegraphics[width=0.49\linewidth]{\figpath{ch05_snapshot_050}}
\caption{標準のprograde--prograde遭遇。時刻は0、0.27、0.35、0.50 Gyrで、全図の軸範囲は$-90$から90 kpc、縦横比は1です。近点後には、初期の円盤から長く伸びる潮汐構造が見えます。色とマーカーは出身銀河に固定しています。Notebookには0.20、0.60 Gyrの図とアニメーションも保存しています。}
\label{fig:ch5-snapshots}
\end{figure}

初めはそれぞれの中心の周りに円盤があります。接近すると円盤が変形し、近点後には外側へ長く伸びる構造と、二つの中心の間をつなぐ構造が現れます。前者を\textbf{潮汐テイル}、後者を\textbf{潮汐ブリッジ}と呼びます。これらは本章の数値結果で粒子が集まって見える構造であり、ガスの衝撃波や新しく生まれた星を描いているのではありません。

\subsection{相手の重力と潮汐力を区別する}
銀河1の周りの星を考えます。相手の銀河2は、星だけでなく銀河1の中心も引きます。円盤の形を変えるのは、全体を同じように並進させる成分ではなく、\textbf{星と銀河中心が受ける加速度の差}です。ゼミで学んだ潮汐力の考え方に対応します。\cite{ryden2025}

本章の標準条件$a_1=a_2=b$では、銀河1の中心加速度は、銀河2のPlummer場を$\Rvec_1$で評価したものと一致します。$\boldsymbol s=\rvec_i-\Rvec_1$とすれば、中心に対する星の相対加速度は
\begin{equation}
\frac{d^2\boldsymbol s}{dt^2}
=\avec_1(\boldsymbol s)
+\underbrace{\left[\avec_2(\Rvec_1+\boldsymbol s)-\avec_2(\Rvec_1)\right]}_{\text{相手の銀河による加速度の差}}
\label{eq:ch5-tidal-difference}
\end{equation}
です。ここで$\avec_2$は銀河2の中心位置も含めて評価する場を表しています。$a_1,a_2,b$を独立に変更した場合には、中心への有効力と中心位置でのPlummer場が一致するとは限らないので、この差の書き方をそのまま流用しません。


% BEGIN TikZ addition: ch05_tidal_difference
図\ref{fig:concept-ch5-tide}では、相手からの加速度そのものと、中心の加速度を引いた差を分けています。遠い側の差が左向きでも、重力が斥力になったわけではありません。この差を、計算済みの両銀河の重力へもう一度足す必要もありません。
\begin{figure}[!htbp]
\centering
% Illustration along the line joining two well-separated galaxies.
% Only the companion's gravitational contribution is drawn.
% Arrow lengths use the point-mass far-field form. This is not the close-passage run.
\begin{tikzpicture}[textbook diagram]
  \node[anchor=west,font=\sffamily\bfseries] at (0,4.05) {(a) Gravity from the companion};
  \node[small label,anchor=west,text=figgray] at (0,3.52) {Companion contribution only; not a simulation snapshot};
  \draw[guide] (0.65,1.72) -- (11.40,1.72);
  \foreach \x/\len/\lab in {2/0.71429/far,3.2/1.04043/c,4.4/1.65406/near} {
    \draw[vector,figblue] (\x,1.72) -- ++(\len,0);
    \draw[guide] (\x,1.77) -- (\x,2.45);
  }
  \fill[figblue] (2,1.72) circle (1.8pt);
  \fill[figblue] (4.4,1.72) circle (1.8pt);
  \draw[black,line width=1pt] (3.12,1.64) -- (3.28,1.80);
  \draw[black,line width=1pt] (3.12,1.80) -- (3.28,1.64);
  \node[above] at (2.0,2.45) {Far side};
  \node[above] at (3.5,2.92) {Center 1};
  \draw[guide] (3.2,2.45) -- (3.5,2.80);
  \node[above] at (5.05,2.45) {Near side};
  \draw[guide] (4.4,2.45) -- (5.05,2.45);
  \node[below] at (2.0,1.72) {$\boldsymbol a_{\rm far}$};
  \node[below,yshift=-5mm] at (3.55,1.72) {$\boldsymbol a_{\rm c}$};
  \node[below] at (5.80,1.72) {$\boldsymbol a_{\rm near}$};
  \draw[figred,dashed] (9.0,1.72) circle (0.60);
  \draw[figred,dashed] (9.0,1.72) circle (0.90);
  \draw[black,line width=1pt] (8.92,1.72) -- (9.08,1.72);
  \draw[black,line width=1pt] (9,1.64) -- (9,1.80);
  \node[align=center] at (9,3.0) {Companion\\(center 2)};
  \node[align=center] at (8.9,0.18) {All three accelerations\\point to the companion.};
  \draw[black!25] (0,-0.50) -- (11.55,-0.50);
  \node[anchor=west,font=\sffamily\bfseries] at (0,-1.10) {(b) Subtract the acceleration of center 1};
  \draw[guide] (0.65,-2.55) -- (11.40,-2.55);
  \fill[figblue] (2,-2.55) circle (1.8pt);
  \fill[figblue] (4.4,-2.55) circle (1.8pt);
  \draw[black,line width=1pt] (3.12,-2.63) -- (3.28,-2.47);
  \draw[black,line width=1pt] (3.12,-2.47) -- (3.28,-2.63);
  \draw[vector,figred] (2,-2.55) -- (1.67386,-2.55);
  \draw[vector,figred] (4.4,-2.55) -- (5.01363,-2.55);
  \node[above] at (1.65,-2.15) {$\boldsymbol a_{\rm far}-\boldsymbol a_{\rm c}$};
  \node[below] at (4.98,-2.73) {$\boldsymbol a_{\rm near}-\boldsymbol a_{\rm c}$};
  \node[below] at (3.2,-2.55) {$0$};
  \draw[figred,dashed] (9,-2.55) circle (0.60);
  \draw[figred,dashed] (9,-2.55) circle (0.90);
  \draw[black,line width=1pt] (8.92,-2.55) -- (9.08,-2.55);
  \draw[black,line width=1pt] (9,-2.63) -- (9,-2.47);
  \node[draw=black!30,fill=white,text width=106mm,align=center,inner sep=5pt] at (5.75,-4.17) {$\delta\boldsymbol a_2(\boldsymbol s)=\boldsymbol a_2(\boldsymbol R_1+\boldsymbol s)-\boldsymbol a_2(\boldsymbol R_1)$};
\end{tikzpicture}

\caption{相手銀河が十分遠い配置での説明図。上段は相手からの加速度、下段は銀河1の中心加速度を引いた差です。銀河1自身の重力は描いていません。$a_1=a_2=b$で中心加速度が対応する場合を考えています。近点8 kpcの標準計算のスナップショットではありません。}
\label{fig:concept-ch5-tide}
\end{figure}
\FloatBarrier
% END TikZ addition: ch05_tidal_difference

相手が十分遠く、$|\boldsymbol s|\ll D$かつ$D\gg a_2$なら、加速度の差の尺度は
\begin{equation}
|\delta\avec_2|\sim\frac{GM_2}{D^3}|\boldsymbol s|
\label{eq:ch5-tidal-scale}
\end{equation}
です。同じ相手までの距離でも、円盤の中心から離れた星ほど差を受けやすいことが分かります。また、円盤外側ほど出身銀河による束縛も弱くなります。ただし、式\eqref{eq:ch5-tidal-scale}だけでテイルの長さや向きが決まるわけではありません。遭遇中の星の回転と、時間に伴って変わる重力を積分した結果を見ます。

\begin{failbox}{潮汐力をもう一度足さない}
式\eqref{eq:ch5-tracer-ode}の二つの重力を使えば、星と中心が受ける重力の差はすでに含まれています。さらに「潮汐力」として別の加速度を足すと二重計算になります。また、本章の近点では円盤外縁15 kpcが中心間距離8 kpcより大きいため、円盤全体に$|\boldsymbol s|\ll D$の展開を使えません。実装は展開式ではなく、両Plummer場を直接評価します。
\end{failbox}

\subsection{出身粒子を拡大図で追跡する}
図\ref{fig:ch5-zoom}は同じ0.35 Gyrの結果を拡大したものです。円盤の内側だけを残して新たに計算したのではありません。中心の間にブリッジがあり、その外側へ曲がったテイルが伸びています。
\begin{figure}[htbp]
\centering
\includegraphics[width=0.76\linewidth]{\figpath{ch05_tidal_zoom}}
\caption{0.35 Gyrの拡大図。軸範囲は$\pm38$ kpcで、図\ref{fig:ch5-snapshots}とは異なります。矢印はブリッジとテイルの位置の例を示します。全体図と同じ粒子・同じ計算結果です。}
\label{fig:ch5-zoom}
\end{figure}

他方の銀河中心の近くに、異なる出身色の星が見えることもあります。しかし、近くへ来たことだけで、その銀河に永久に捕獲されたとは判断できません。ここではまず出身ラベルを保ったまま、どこへ移動したかを確認します。複雑な束縛判定やテイルの自動抽出は、本章には追加しません。

\FloatBarrier
\section{外へ移動した星を数える}
\label{sec:ch5-outer}
\subsection{何からの距離を測るかを決める}
図の見た目を補うため、出身銀河の中心から遠くへ移った粒子の割合を測ります。粒子$i$の出身銀河を$k(i)$と書くと、その距離は
\begin{equation}
s_i(t)=|\rvec_i(t)-\Rvec_{k(i)}(t)|
\label{eq:ch5-origin-radius}
\end{equation}
です。原点や、その時刻に最も近い銀河中心からの距離ではありません。

本章では閾値を$s_{\rm cut}=30\unit{kpc}=2R_{\max}$と定めます。各出身群$k$について
\begin{equation}
f_k(t)=\frac{1}{N_k}\sum_{i:\,k(i)=k}\boldsymbol{1}\left(s_i(t)>s_{\rm cut}\right)
\label{eq:ch5-outer-fraction}
\end{equation}
を計算します。$\boldsymbol{1}$は条件が真なら1、偽なら0です。初期配置の範囲は15 kpc以内なので、初期値では$f_1=f_2=0$です。
\inputcode{chap05/outer-fraction}

\texttt{history[:, origin, :2]}は、各テスト粒子の出身中心の位置を各時刻で取り出します。\texttt{history[:, 2:, :2]}と同じ$(N_{\rm save},N,2)$の形になるので、その差から距離を測れます。粒子を除去して分母を変える操作はしません。

\subsection{数値の意味と限界を確認する}
標準計算の0.60 Gyrでは、$f_1=\chfivevalOuterOne$、$f_2=\chfivevalOuterTwo$でした。1000粒子ずつの標本なので、各割合は0.001刻みです。二つは等質量・同じ分布モデルですが、有限の乱数標本が違うため、数値は完全には一致しません。
\begin{table}[htbp]
\centering
\caption{出身中心から30 kpcを越えた粒子の割合。これは位置による指標で、非束縛粒子や永久に剥ぎ取られた質量の割合ではありません。}
\label{tab:ch5-fractions}
\begin{tabular}{rrr}\toprule
$t\ [\mathrm{Gyr}]$ & $f_1$ & $f_2$ \\
\midrule
0.00 & 0.000 & 0.000 \\
0.20 & 0.000 & 0.000 \\
0.27 & 0.000 & 0.000 \\
0.35 & 0.023 & 0.017 \\
0.50 & 0.404 & 0.413 \\
0.60 & 0.425 & 0.432 \\
\bottomrule\end{tabular}

\end{table}
\begin{figure}[htbp]
\centering
\includegraphics[width=0.76\linewidth]{\figpath{ch05_outer_fraction}}
\caption{式\eqref{eq:ch5-outer-fraction}の時間変化。青実線と朱色破線が各出身群です。中心間距離が増えても、星の軌道はすぐに初期の円運動へ戻るわけではありません。}
\label{fig:ch5-fractions}
\end{figure}

この指標は\textbf{「外へ移動した割合」}です。出身銀河に束縛されたまま大きな軌道を回る粒子、後で戻る粒子、他の銀河の近くへ移った粒子も含み得ます。閾値30 kpcは本章で定めた測定条件であり、自然に一意に決まる潮汐半径ではありません。「剥ぎ取りが起こり得る様子」を見ることと、永久に失われた質量を定量的に予測することを区別します。

\FloatBarrier
\section{星のエネルギーが変化する理由を考える}
\label{sec:ch5-energy}
第4章では固定した一つのPlummer場を使ったため、各星の比エネルギーが保存されました。今回は両中心が動きます。慣性系での比エネルギーを
\begin{equation}
\varepsilon_i(t)=\frac12|\vvec_i(t)|^2+\Phi_1(\rvec_i,t)+\Phi_2(\rvec_i,t)
\label{eq:ch5-tracer-energy}
\end{equation}
と定義しても、一般には一定ではありません。

運動方程式を使って時間微分すると、$\vvec_i\cdot\avec_i$と$\vvec_i\cdot\nabla\Phi$が打ち消し合い、
\begin{equation}
\frac{d\varepsilon_i}{dt}=\frac{\partial\Phi}{\partial t}
=\sum_{k=1}^2\Vvec_k\cdot\avec_{k,i}
\label{eq:ch5-energy-rate}
\end{equation}
が残ります。最後の等号は、$\Phi_k$が$\rvec_i-\Rvec_k(t)$の関数で、$M_k,a_k$は一定であることから得られます。これは本章のモデルの式を微分した関係です。\textbf{動力学的摩擦を入れなくても、動く重力場は星のエネルギーを変えます。}
\inputcode{chap05/tracer-energy}

図\ref{fig:ch5-tracer-energy}では、銀河1の出身粒子のうち、初期半径が$3R_d=9$ kpcに最も近い1粒子を選びました。その初期半径は$\chfivevalWatchRadius\unit{kpc}$です。これは個別の軌道の例で、星全体の平均値ではありません。
\begin{figure}[htbp]
\centering
\includegraphics[width=0.76\linewidth]{\figpath{ch05_tracer_energy}}
\caption{選んだ1粒子の比エネルギー変化。規格化は$GM_1/a_1$です。図\ref{fig:ch5-center-energy}の「中心の保存量の誤差」と違い、こちらの変化には時間依存する場の物理的な効果が含まれます。}
\label{fig:ch5-tracer-energy}
\end{figure}

星へのトルクも一般には0ではないので、慣性系の原点の周りの各星の角運動量も保存しません。また、星から中心へ反作用を返していないため、星の表示上の重みにエネルギーを掛け、その和を$E_{\rm c}$へ足して「全系の保存量」とすることはできません。今回の確認では、\textbf{中心の二体系の保存量}と、\textbf{星の位置の時間刻み依存性}を分けます。

\section{相互作用計算の時間刻みを確認する}
\label{sec:ch5-convergence}
\subsection{初期粒子も評価時刻も固定する}
二銀河が動く間の星には、第4章の孤立円盤の円軌道の解析位置は使えません。その位置からずれた量には、意図した相互作用の効果が入るためです。ここでは第3章の非円軌道と同じく、\textbf{同じモデルの、異なる刻みの数値解どうし}を比較します。

初期配列、終時刻0.60 Gyr、評価する151個の時刻を固定します。ステップ数だけを3000、6000、12000、24000とし、刻みは0.20、0.10、0.05、0.025 Myrです。比較時刻を4 Myrごとにそろえるため、各計算で\texttt{save\_every = n\_steps // 150}とします。この保存設定は、標準図の301時刻とは別の検証用設定です。

\begin{table}[htbp]
\centering
\caption{刻みごとの中心軌道と0.60 Gyrの外側粒子割合。$D_{\min}$と$e_{E,{\rm c}}$は中心の全ステップから求めています。$f_1,f_2$の閾値は30 kpcです。}
\label{tab:ch5-resolution}
\small\begin{tabular}{rrrrrr}\toprule
$N_{\rm step}$ & $\Delta t\ [\mathrm{Myr}]$ & $D_{\min}\ [\mathrm{kpc}]$ & $e_{E,{\rm c}}$ & $f_1$ & $f_2$ \\
\midrule
3000 & 0.200 & 8.000109881 & $5.28\times 10^{-5}$ & 0.425 & 0.432 \\
6000 & 0.100 & 8.000014569 & $1.32\times 10^{-5}$ & 0.425 & 0.432 \\
12000 & 0.050 & 8.000004531 & $3.30\times 10^{-6}$ & 0.425 & 0.432 \\
24000 & 0.025 & 8.000001962 & $8.26\times 10^{-7}$ & 0.425 & 0.432 \\
\bottomrule\end{tabular}

\end{table}

中心の最小距離は8 kpcへ近づき、エネルギー誤差も刻みを細かくするにつれて減ります。一方、この表では外側粒子の割合が全行で同じ値でした。しかし、その一致だけで各星の位置まで正確だとは判断できません。割合は粒子が閾値のどちら側にいるかだけを数え、個々の位置の違いを捨てる量だからです。

\subsection{隣接する刻みの位置の差を測る}
粗い刻みを$\Delta t$、細かい刻みを$\Delta t/2$と書きます。評価する各時刻$t_j$で位置の差を取り、中心と星について
\begin{align}
\delta_{\rm c}(\Delta t)&=\max_{j,k}\frac{|\Rvec_{k,\Delta t}(t_j)-\Rvec_{k,\Delta t/2}(t_j)|}{D_0},\notag\\
\delta_{\rm t}(\Delta t)&=\max_{j,i}\frac{|\rvec_{i,\Delta t}(t_j)-\rvec_{i,\Delta t/2}(t_j)|}{R_d}
\label{eq:ch5-self-difference}
\end{align}
を求めます。これは真の解との誤差ではなく、二つの数値解の差です。同じ時刻・同じ粒子番号で比較します。粒子を位置順に並べ替えてから対応させることはしません。

少数の軌道に大きな差が集中する場合を見分けるため、星については二乗平均平方根の差も求めます。
\begin{equation}
\delta_{\rm rms}(\Delta t)=\frac{1}{R_d}
\sqrt{\frac{1}{N_{\rm eval}N}\sum_{j,i}|\rvec_{i,\Delta t}(t_j)-\rvec_{i,\Delta t/2}(t_j)|^2}.
\label{eq:ch5-rms-difference}
\end{equation}
$N_{\rm eval}=151$、$N=2000$です。最大値と平均的な大きさの両方を残し、一方だけで精度を説明しないようにします。
\inputcode{chap05/time-convergence}

\begin{table}[htbp]
\centering
\caption{151個の共通時刻での位置の自己収束。左の2列はステップ数です。前行との比は$\delta_{\rm t}$から求めています。$\delta_{\rm rms}$は星・評価時刻の両方にわたる二乗平均平方根です。}
\label{tab:ch5-self-convergence}
\small\begin{tabular}{rrrrrr}\toprule
粗い計算 & 細かい計算 & $\delta_{\rm c}$ & $\delta_{\rm t}$ & $\delta_{\rm rms}$ & 前行との比 \\
\midrule
3000 & 6000 & $3.90\times 10^{-6}$ & $2.51\times 10^{-1}$ & $2.68\times 10^{-3}$ & --- \\
6000 & 12000 & $9.75\times 10^{-7}$ & $6.41\times 10^{-2}$ & $6.71\times 10^{-4}$ & 3.92252 \\
12000 & 24000 & $2.44\times 10^{-7}$ & $1.61\times 10^{-2}$ & $1.68\times 10^{-4}$ & 3.97963 \\
\bottomrule\end{tabular}

\end{table}

最も細かい比較では、最大差$\delta_{\rm t}=\chfivevalLastDifference$、二乗平均平方根は$\delta_{\rm rms}=\chfivevalLastRMS$でした。$R_d=3$ kpcなので、最大の位置差はおよそ$\chfivevalLastDifferenceKpc$ kpcです。二つの値の違いは、一部の粒子が全体の平均より刻みに敏感であることを示します。

\begin{figure}[htbp]
\centering
\includegraphics[width=0.76\linewidth]{\figpath{ch05_convergence}}
\caption{隣接する刻みの星の位置の差。青丸が最大値、朱色四角が二乗平均平方根です。灰色破線は最も細かい比較の最大値を通る傾き2の参照線で、実測点への当てはめではありません。}
\label{fig:ch5-convergence}
\end{figure}

解の主な誤差が$C(\Delta t)^2$で表せる範囲なら、隣接する解の差は$(3/4)C(\Delta t)^2$です。その差どうしの比は約4になります。表\ref{tab:ch5-self-convergence}の細かい側では比が\chfivevalLastRatio となり、4へ近づいています。標準の0.05 Myrは、この比較で潮汐形態と外側割合を調べるための基準とします。すべての粒子の真の軌道誤差が一定値以下だと保証する設定ではありません。より長い計算、より小さい近点、別の密度集中度へ変更したら、刻みも再確認します。

\subsection{一致する数値の意味を確かめる}
自己収束だけでは、加速度の式そのものが正しいことを証明できません。中心間力の作用・反作用、星への重力の和、中心軌道だけの独立計算、単独Plummer場中の円盤の検証を組み合わせます。これらの独立な確認コードも配布物に含めています。

全中心・全星に同じ初期位置$(5,-3)$ kpcと速度$(30,-20)$ km/sを加える試験も行います。相対位置は変わらないので、同じ相対運動になるはずです。今回の出身中心に対する粒子位置の差は、$R_d$で規格化した最大値で$\chfivevalBoostDiff$でした。標準の刻みを変えた差より十分小さくなっています。ただし、軌道が摂動に敏感な場合には丸めの違いも時間とともに増え得るので、任意の計算で同じ数値になるとは考えません。

\FloatBarrier
\section{このモデルで分かることと限界を整理する}
\label{sec:ch5-limits}
本章の計算から、二つの銀河が接近すると、星の軌道が変わり、円盤からテイルやブリッジが生じる様子を確認できました。新しい星を生成したわけではなく、初めからあった星を別の位置へ運んだ結果です。また、星のエネルギーが変化しても、中心の二体系の軌道エネルギーは保存されるモデルだと確認しました。

中心が再び離れていくことは計算の失敗ではありません。$\varepsilon=0$の中心軌道に摩擦や星からの反作用がないので、近接通過後に自動的に合体へ切り替わる仕組みはありません。ポテンシャルの形も質量も固定なので、潮汐変形を受けて銀河の重力場自体が変わる過程、質量損失に伴う中心軌道の変化、合体残骸が落ち着く過程は扱っていません。

また、円盤と遭遇軌道は同じ平面に置いています。円盤の傾斜角、厚み、鉛直方向の運動は扱いません。したがって、後のprograde--retrogradeの比較は共面配置で内部回転の向きだけを変える実験です。任意の傾斜を持つ銀河遭遇の形態まで、この二つの向きだけで分類することはできません。

次の第6章では、完成したコードを用い、近点距離、質量比、円盤の回転方向を一つずつ変えます。さらに発展では、衛星銀河の中心へChandrasekhar動力学的摩擦を加え、軌道減衰を調べます。その速度依存力にはRK4を使います。摩擦を質量ゼロの星一つずつへ掛けるわけではありません。ここまでの重力だけの計算を基準として、何を追加したかを明確にします。

\section{本章のまとめ}
本章では、2中心と多数のテスト粒子を一つの配列へまとめました。中心同士は共通の有効対ポテンシャルで動かし、星には二つの移動するPlummer場を足しました。中心間の有効力と、広がった密度分布間の厳密な相互力は区別しました。

初期軌道は、採用したポテンシャルの比エネルギーと比角運動量から求めました。各円盤は第4章の方法で作り、位置と速度の両方に中心の運動を加えました。KDKでは全中心・全粒子をそろえて更新し、新しい中心位置を使って星の加速度を求め直しました。

結果には潮汐テイルとブリッジが現れました。これらを重力の差による軌道の変化として読み、出身中心から遠くへ移動した星の割合も測りました。ただし、その割合を非束縛率や質量損失率とは呼びません。中心の保存量、星の位置の自己収束、異なる座標系での比較を使い、図の見た目だけで検証を終えないことを確認しました。

\chapterexercises
\paragraph{演習5.1　確認：どの重力を計算するか}
中心1、中心2、銀河1出身の星、銀河2出身の星の加速度をそれぞれ書いてください。$M_1\ne M_2$、$a_1\ne a_2$の場合でも、共通の$b$を使った中心間力が作用・反作用を満たすことを示してください。\texttt{origin}を星の重力源の選別に使ってはいけない理由も述べてください。

\paragraph{演習5.2　確認：有効軌道の初期条件}
表\ref{tab:ch5-standard}の値から、$h,U_t,U_r$を求め、$\varepsilon(0)=0$を確認してください。同じ$h$で$D=D_p$としたとき、動径速度が0になることも確かめてください。$b\to0$の極限で、第2章の放物線の式へ戻ることを示してください。

\paragraph{演習5.3　確認：円盤の並進}
\texttt{place\_disk}の出力から中心位置と中心速度を引き、元の局所円盤へ戻ることを確認してください。中心速度を足し忘れた場合、各粒子の中心に対する速度にどの項が加わるでしょうか。その崩れを、正しい潮汐テイルの再現と判断できるでしょうか。

\paragraph{演習5.4　応用：時間刻みと誤差指標}
同じ初期配列で表\ref{tab:ch5-resolution}と表\ref{tab:ch5-self-convergence}を再計算してください。外側割合が同じ値でも最大位置差が0ではない理由を説明してください。自己収束の差を真の軌道誤差と区別し、標準設定を変更した場合に再確認すべきものを挙げてください。

\paragraph{演習5.5　応用：保存間隔だけを変える}
標準の時間刻みを固定します。\texttt{save\_every}を1、40、200として比べてください。全粒子の最終状態、毎ステップの中心履歴、全履歴の容量はどう変わりますか。出力の細かさと時間積分の精度を区別してください。メモリが少ない場合は、粒子数を減らしてから三つの保存設定を比較して構いません。

\paragraph{演習5.6　応用：系全体を等速で動かす}
すべての中心と星に、同じ初期位置$(5,-3)$ kpcと速度$(30,-20)$ km/sを加えてください。出身中心からの相対位置と中心間距離を、もとの計算と比べてください。慣性系での各星のエネルギーの数値まで、もとの計算と一致する必要があるでしょうか。

\paragraph{演習5.7　発展：片方の円盤を逆回転にする}
標準の初期位置と中心の初期速度を保ち、銀河2出身の星の\textbf{中心に対する速度だけ}を反転してください。0.60 Gyrの粒子分布と$f_1,f_2$を比べてください。星の速度全体へ単に$-1$を掛けた場合とは何が違いますか。中心軌道と銀河1出身の星の軌道が変わらないことも確認してください。この比較を、自己重力銀河全体の逆回転合体の定量的予測とみなしてよいでしょうか。
